% !TeX root = ../../main.tex

\section{Clinical Motivation}

\subsection{Epidemiology of Kidney Transplantations and Failures Thereof}

% Paragraph 1: ESKD is a growing public health problem
%   - ESKD is the final stage of CKD (both are loosely defined)
%   - Risk factors for ESKD (diabetes, hypertension, glomerulonephritis etc.)
%   - Incidence (# new cases), prevalence (total disease burden), and mortality of ESKD. In 2 sentences or less and/or as a table, give it for:
%     - Globally
%     - OECD nations (developed nations)
%     - United States
%     - North Carolina, specifically
%   - Economic burden of ESKD (globally, and per capita in US)
% - Quality of life implications for patients with ESKD

End-stage kidney disease (ESKD), also known as kidney failure, represents a significant and growing public health challenge worldwide~\cite{bikbov_Global_2020}. ESKD is the final stage of chronic kidney disease (CKD), characterized by the irreversible loss of kidney function to the point where homeostasis and survival cannot be sustained~\cite{shringi_EndStage_2026}. Though its definition and surveillance techniques greatly vary by country, ESKD is a globally relevant condition that is thought to affect at least 2 million people to the point where kidney replacement therapies are required~\cite{trillini_Epidemiology_2017,gupta_Epidemiology_2021}. Notably, this number may be severely underestimated since international epidemiologic surveys tend to exclude ESKD patients who do not receive treatment~\cite{thurlow_Global_2021}, a factor that disproportionally impacts ESKD patients in the Global South. However, in higher-income countries where CKD treatments are more prevalent, ESKD incidence is higher where risk factors such as diabetes, hypertension and other cardiovascular diseases, obesity, and greater healthcare spending are more common~\cite{thurlow_Global_2021}. This, of course, implicates the United States: ESKD affects approximately 14.9 \% of the population (over 50 million individuals), with 5-year survival rates from ESKD onset ranging from 41 \% to 47 \% depending on the type of dialysis administered~\cite{gupta_Epidemiology_2021}.

It is critical to recognize that death is not the only adverse outcome of ESKD; many patients live with significant morbidity and diminished quality of life~\cite{nayana_cross_2017}. ESKD patients experience a range of debilitating symptoms, including fatigue, cognitive impairment, and cardiovascular complications - all of which collectively impair daily functioning, occupational productivity, social capacity, and overall well-being. The macroeconomic burden of ESKD is therefore substantial, with global per capita expenditures of ESKD treatments reaching between \$19,000 and \$27,000 annually, noted by the United States' National Institute of Diabetes and Digestive and Kidney Health (NIDDK) to be "an immense sum when the per capita [gross domestic product] worldwide (adjusted for purchasing power parity) is approximately \$22,500"~\cite{unitedstatesrenaldatasystem_USRDS_2025}. This risk for financial hardship is still significant in the United States to the point that ESKD is the only disease state for which the U.S. federal government provides universal healthcare coverage, regardless of age, socioeconomic status, or veteran status~\cite{centerformedicareandmedicaidservices_EndStage_2024}. While this policy underscores the critical importance of addressing ESKD, it also highlights the immense strain that this condition places on healthcare systems and resources~\cite{unitedstatesrenaldatasystem_USRDS_2025,gupta_Epidemiology_2021}.

% Paragraph 2: Kidney transplantation is definitive treatment for ESKD
%   - Advantages of kidney transplantation over dialysis
%   - Survival rates of kidney transplantation
%     - Globally
%     - OECD nations (developed nations)
%     - United States
%     - North Carolina, specifically
%   - Economic benefits of kidney transplantation
%   - Quality of life improvements for patients with kidney transplantation

Kidney transplantation\footnote{
    While other sources of organ transplants exist (e.g. \cite{artilesmedina_Kidney_2022} for a case report on kidney autografts, \cite{steadman_United_2025} for a preclinical demonstration of renal xenotransplantation), these alternative approaches are not known to be standard clinical practice. Therefore, this dissertation will treat the term "transplants" as a synonym to "allografts".
} stands as the definitive treatment for ESKD, offering significant advantages over dialysis in terms of survival, quality of life, and economic burden. Transplant recipients generally experience improved survival rates compared to those on dialysis; 1-year survival rates for graft recipients being 12.9 \% higher (85.1 \% for dialysis patients, 98.0 \% for graft recipients)~\cite{unitedstatesrenaldatasystem_USRDS_2025}, with the gap widening over longer timespans. Medicare recipients in the United States also see a financial benefit to receiving a kidney transplant, as their per-person-per-year fee-for-service costs drop from \$103,000 as of 2023 inflation levels for in-center hemodialysis to just under \$45,000 - a 56.7 \% reduction~\cite{unitedstatesrenaldatasystem_USRDS_2025}. Furthermore, kidney transplantation markedly enhances patients' quality of life by restoring renal function, alleviating the physical and psychological burdens associated with dialysis, and enabling greater social and occupational engagement.

% Paragraph 3: Kidney transplant rejection remains a significant clinical challenge
%   - Acute and hyper-acute rejection is adequately addressed
%   - Chronic rejection is remaining challenge
%   - Statistics on graft survival rates
%   - Graft survival becomes worse with inflammation and fibrosis

However, kidney transplantation is not without its challenges; transplant rejection remains a significant clinical hurdle that compromises long-term graft survival. While acute and hyper-acute rejection episodes are generally well-managed with current immunosuppressive therapies, chronic rejection continues to pose a formidable obstacle~\cite{poggio_Longterm_2021}. The 10-year renal graft survival rate in the United States is 53.6 \%~\cite{hariharan_LongTerm_2021}, much lower than the 1- and 5-year rates previously mentioned. At this time scale, recipient death, graft rejections, age, and the presence of donor serum antibodies (DSA) are the primary factors associated with graft loss~\cite{unitedstatesrenaldatasystem_USRDS_2025,betjes_Causes_2022,el-zoghby_Identifying_2008,kanbay_metaanalysis_2025}. The presence of inflammation and fibrosis within the renal allograft further exacerbates this issue, as these pathological changes are closely associated with declining graft function and eventual failure~\cite{haas_Differences_2017,betjes_Causes_2022,gago_Kidney_2012,nankivell_Rejection_2010,saritas_Kidney_2021}.

\subsection{Overview of Relevant Pathologies in Kidney Allografts}

% Paragraph 1: what really is graft rejection and loss, and why does it happen?
%   - Graft rejection = autoimmune response against donor kidney
%   - Graft loss = irreversible dysfunction of donor kidney
%   - Not mutually exclusive; graft rejection can be halted / graft loss can occur without rejection

Fundamentally, graft rejection is the autoimmune reaction of a transplant recipient against their donor organ. This foreign body response can manifest in various forms, including T cell-mediated rejection (TCMR) and antibody-mediated rejection (AMR), both of which can occur acutely or chronically~\cite{nankivell_Rejection_2010,loupy_Banff_2020,schinstock_Banff_2019}. Rejection does not invariably lead to graft loss; with timely and appropriate intervention, rejection episodes can often be halted, preserving graft function~\cite{callemeyn_Revisiting_2021,callemeyn_Allorecognition_2022}. However, rejection remains a primary contributor to graft dysfunction as well as a risk factor for eventual graft loss~\cite{poggio_Longterm_2021,hariharan_LongTerm_2021}.
  
% Paragraph 2: graft dysfunction and loss can happen due to multiple reasons:
%   - Donor-related factors
%     - ABO incompatibility, other immunologic mismatches
%     - Pre-existing conditions (Hep C, HIV, diabetes, hypertension etc.)
%   - Transplantation-related factors
%     - Ischemia-reperfusion injury
%     - Surgical complications
%     - Delayed graft function
%     - Hyperacute rejection
%   - Environmental factors
%     - Medications (nephrotoxic drugs etc.)
%     - Lifestyle factors (diet, exercise, smoking etc.)
%   - Recipient-related factors
%     - Comorbidities (diabetes, hypertension etc.)
%     - Immunosuppressant adherence
%     - Transient immune activity (infections, trauma etc.)
%     - Rejection

Graft dysfunction and loss can arise from a multitude of factors spanning donor characteristics, transplantation procedures, environmental influences, and recipient-specific variables~\cite{alexandre_AntibodyMediated_2018}. Donor-related factors are characterized by risk factors presented in the donor organ or the donor's medical history. This includes serotype incompatibilities (e.g., ABO blood group mismatches, HLA mismatches) as well as pre-existing conditions such as hepatitis C, HIV, diabetes mellitus, and hypertension~\cite{shringi_EndStage_2026,loupy_Banff_2020}. Transplantation-related factors encompass complications arising during or immediately after the surgical procedure. This includes ischemia-reperfusion injury, surgical complications, delayed graft function, and hyperacute rejection. Environmental factors refer to external influences that may impact graft health, such as exposure to nephrotoxic medications and lifestyle choices including diet, exercise, and smoking habits. Finally, recipient-related factors involve individual patient characteristics and behaviors. These factors include comorbidities such as diabetes and hypertension, adherence to immunosuppressive regimens, transient immune activations due to infections or trauma, and episodes of rejection.

% Paragraph 3: rejection is the leading cause of graft dysfunction and loss
%   - Two main, overlapping types of rejection:
%     - T cell-mediated rejection (TCMR)
%     - Antibody-mediated rejection (AMR)
%   - TCMR and AMR can co-occur, be acute or chronic
%   - Epidemiology of rejection
%     - Incidence rates
%     - Timing post-transplant
%     - Impact on graft survival

Out of these various factors, rejection remains the leading cause of graft dysfunction and loss~\cite{hariharan_LongTerm_2021}. TCMR and AMR represent the two primary, overlapping types of rejection that can afflict kidney allografts. Both TCMR and AMR can be classified into active and chronic forms (formerly denoted as "acute" and "chronic" rejection, respectively)~\cite{loupy_Subclinical_2015,loupy_Banff_2020}, with the former characterized by rapid onset and the latter by gradual progression. Since TCMR and AMR in the kidney are not mutually exclusive pathologies, both types of rejection can exacerbate the risk of graft loss within the first 10 years of transplantation~\cite{matignon_Concurrent_2012,betjes_Causes_2022}. Betjes \emph{et al.} found that TCMR and AMR together constituted the leading known, death-censored causes for graft loss, accounting for 58 \% of such cases in their cohort~\cite{betjes_Causes_2022}.

\subsection{Limitations of Current Clinical Practices}

% Paragraph 1: lab-based biomarkers only assess downstream effects of pathology
%   - Serum creatinine and proteinuria are nonspecific, lagging
%   - eGFR (and "accurate" alternatives) are unreliable

To monitor the health of kidney allografts, clinicians currently rely on a combination of laboratory-based biomarkers, histopathological evaluations, and imaging modalities. The simplest and most commonly used laboratory biomarkers include serum creatinine levels and proteinuria measurements~\cite{goutaudier_Evaluation_2024,naesens_Proteinuria_2016}. However, these markers are inherently nonspecific and often lag behind the actual pathological changes occurring within the graft. Likewise, proteinuria can be influenced by a variety of non-rejection-related factors such as glomerular diseases and polyomavirus infections~\cite{naesens_Proteinuria_2016}. Other metrics such as estimated glomerular filtration rate (eGFR) may also be measured, though they also suffer from reliability issues due to confounding variables like BMI, hydration status, and certain medications~\cite{loupy_Subclinical_2015,sasaki_Variability_2025}.

% Paragraph 2: biopsies are the gold standard, but invasive, risky, and late

As a counterweight to the limitations of laboratory biomarkers, renal allograft biopsies provide a specific, direct assessment of graft pathology, enabling them to act as the gold standard for diagnosing rejection. Biopsies involve the photography and analysis of tissue samples obtained via percutaneous needle insertion into the graft. The small tissue samples are then stained and examined under a microscope by a trained pathologist, who looks for histological features indicative of rejection, inflammation, fibrosis, and other pathological changes~\cite{loupy_Banff_2020}. However, biopsies are invasive procedures that carry inherent risks such as bleeding, infection, and graft injury. Additionally, biopsies are typically performed only when clinically indicated (e.g., significant rise in serum creatinine), which may be too late to prevent irreversible graft damage~\cite{schinstock_Value_2017}. Furthermore, sampling errors and inter-observer variability among pathologists can lead to inconsistent diagnoses~\cite{schinstock_Banff_2019}.

% Paragraph 3: multi-omic approaches are emerging but expensive
%   - Donor-derived cfDNA
%   - transcriptomics?

The practical limitations of biopsies have spurred interest in emerging multiomic approaches that aim to provide a more comprehensive and noninvasive assessment of graft health. One such approach involves the detection of donor-derived cell-free DNA (dd-cfDNA) in the recipient's bloodstream~\cite{bloom_CellFree_2017}. Elevated levels of dd-cfDNA have been associated with AMR~\cite{bloom_CellFree_2017} as well as cases of subclinical rejection~\cite{aubert_Cellfree_2024}, offering a promising biomarker for early detection. However, the implementation of dd-cfDNA testing is hindered by a lack of methodological standardization, sensitivity issues, high costs, and the need for specialized laboratory infrastructure~\cite{chua_Challenges_2025,song_Limitations_2022}. Other multiomic techniques, such as the mRNA, miRNA, and proteomic profiling of samples of peripheral blood, urine, or even the gut microbiome, are also being explored for their potential to identify molecular signatures of rejection~\cite{goutaudier_Evaluation_2024,marquet_thorough_2025,vanloon_Diagnostic_2021,lubetzky_Urinary_2021,gielis_Combined_2021,holle_Gut_2025}. While these approaches hold promise, they too face challenges related to cost, complexity, and the need for further validation in clinical settings.

% Paragraph 4: imaging could help - but limited
%   - Scintigraphy: it's just eGFR with more steps
%   - MR elastography: expensive and impractical
%   - Ultrasound B-mode is too nonspecific, confounded
%   - Ultrasound/Doppler only provides part of the picture

Another approach to monitoring renal allograft health involves medical imaging modalities, which can provide noninvasive visualization of the graft's structure and function. The most analogous technique to the previously described laboratory biomarkers is scintigraphy, where radioactive tracers are injected into the bloodstream and tracked such that their uptake and excretion by the graft can be visualized~\cite{alnazer_Recent_2021}. While this approach has allowed for the direct monitoring of graft function, it essentially mirrors eGFR assessments with added complexity, radiation exposure, cost, and replication challenges. Magnetic resonance elastography (MRE) offers a noninvasive means of assessing tissue stiffness by applying mechanical vibrations and measuring resulting tissue displacements using magnetic resonance imaging (MRI) techniques~\cite{lee_MR_2012}. Although MRE has shown promise in evaluating fibrosis in kidney allografts, it is also expensive, time-consuming, and inconsistent in performance~\cite{wilson_More_2020} These practical limitations hinder the widespread adoption of MRE for routine graft monitoring.

Ultrasound imaging presents a more accessible and cost-effective option for graft monitoring~\cite{zhou_Current_2021}. Conventional "B-mode" (brightness mode) ultrasound provides real-time images of anatomical structures within the graft, allowing for the rapid assessment of size, shape, and gross abnormalities. This approach is widely used due to its accessibility and safety, especially since even simple metrics such as kidney length have been correlated with graft outcomes~\cite{carnahan_SolidOrgan_2024}. However, many pathological changes within the graft (e.g., inflammation, fibrosis) do not manifest as overt anatomical alterations, rendering B-mode ultrasound insufficiently sensitive and specific for early detection of graft dysfunction.

However, in a similar manner to MRI and MRE, ultrasound imaging can also be used to assess more than just anatomical information. Doppler ultrasound techniques are a common example where blood flow magnitude and direction are estimated within the graft vasculature~\cite{rifkin_Evaluation_1987,irshad_review_2008,sugi_Imaging_2019}.\footnote{
    Despite the name's implication, Doppler ultrasound does not rely on demodulating frequency shifts in ultrasonic echoes. Rather, it tracks tissue motion over time and estimates the direction and speed at which that occurs.
} Likewise, contrast-enhanced ultrasound (CEUS) can also be employed to facilitate the visualization of renal bloodflow and perfusion~\cite{zhou_Current_2021}. Since changes in blood flow such as arterial stenosis and venous thrombosis can significantly impact graft health, Doppler and CEUS provide valuable functional information that complements B-mode imaging. Nevertheless, these subfields of ultrasonography only provides a partial picture of graft health, as it primarily focuses on functional and vascular parameters; inflammation and fibrosis also manifest beyond just blood flow alterations~\cite{saritas_Kidney_2021}.

% Motivating statement: there's a need for a tool that noninvasively detects renal transplant rejection earlier than clinically indicated biopsies while providing relevant information about tissue microstructure

In summary, current clinical practices for monitoring renal allografts are fraught with limitations that hinder early detection and effective management of graft dysfunction and rejection. Laboratory biomarkers are nonspecific and lagging, biopsies are invasive and risky, multiomic approaches are expensive and complex, and existing imaging modalities provide only partial insights into graft health. Therefore, there exists a critical need for a noninvasive tool that can detect kidney transplant rejection earlier than clinically indicated biopsies while providing relevant information about tissue microstructure. Such a tool would significantly enhance the ability to monitor graft health, enabling timely interventions that could improve long-term outcomes for transplant recipients.
